\documentclass[a4paper,11pt]{article}

% ─── Geometry ────────────────────────────────────────────────────────────────
\usepackage{geometry}
\geometry{margin=1in, headheight=15pt}

% ─── Mathematics ─────────────────────────────────────────────────────────────
\usepackage{amsmath}
\usepackage{amssymb}
\usepackage{mathtools}

% ─── Graphics / Colours ──────────────────────────────────────────────────────
\usepackage{graphicx}
\usepackage{xcolor}
\definecolor{blue3}{RGB}{0,0,180}
\definecolor{green3}{RGB}{0,120,0}
\definecolor{orange3}{RGB}{200,100,0}
\definecolor{red3}{RGB}{180,0,0}
\definecolor{grey3}{RGB}{80,80,80}

% ─── Units ───────────────────────────────────────────────────────────────────
\usepackage{siunitx}
\sisetup{detect-all}

% ─── Links (hyperref BEFORE cleveref) ────────────────────────────────────────
\usepackage[colorlinks=true,linkcolor=blue3,citecolor=blue3,urlcolor=blue3]{hyperref}
\usepackage{cleveref}

% ─── Boxes ───────────────────────────────────────────────────────────────────
\usepackage{tcolorbox}
\tcbuselibrary{skins,breakable}

\newtcolorbox{hintbox}[1][]{
    colback=orange3!8,
    colframe=orange3,
    title={Hint:},
    fonttitle=\bfseries,
    #1
}

% ─── Tables / Lists ──────────────────────────────────────────────────────────
\usepackage{booktabs}
\usepackage{enumitem}

% ─── Notation commands ───────────────────────────────────────────────────────
\newcommand{\ktilde}{\tilde{k}}
\newcommand{\ytilde}{\tilde{y}}
\newcommand{\ctilde}{\tilde{c}}
\newcommand{\htilde}{\tilde{h}}
\newcommand{\dd}{\mathrm{d}}
\newcommand{\bgp}{\mathrm{BGP}}

% ─── Header / Footer ─────────────────────────────────────────────────────────
\usepackage{fancyhdr}
\pagestyle{fancy}
\fancyhf{}
\lhead{\textbf{14E022 Macroeconomics I} \quad Practice Exam 2}
\rhead{BSE 2025--26}
\cfoot{\thepage}

% ═════════════════════════════════════════════════════════════════════════════
\begin{document}
% ═════════════════════════════════════════════════════════════════════════════

\begin{center}
    {\LARGE \textbf{14E022 Macroeconomics I}}\\[6pt]
    {\Large \textbf{Practice Exam 2}}\\[8pt]
    {\large BSE --- Academic Year 2025--26}
\end{center}

\vspace{8pt}
\noindent\textit{Instructions: the exam is made of 3 problems and a set of short questions
totalling \textbf{100 points}. You have \textbf{2 hours}. Show all derivations clearly.
For short questions, concise and precise answers score highest.}
\vspace{4pt}

\noindent\rule{\textwidth}{0.4pt}

% ─────────────────────────────────────────────────────────────────────────────
\section{Question 1 --- Augmented Solow Model with Human Capital
\hfill (25 points)}
% ─────────────────────────────────────────────────────────────────────────────

Consider a Solow model augmented with human capital \`a la Mankiw--Romer--Weil (1992).
Output is produced according to
\[
    Y_t = K_t^{\alpha}\,H_t^{\beta}\,(A_t L_t)^{1-\alpha-\beta},
    \qquad \alpha,\beta > 0,\quad \alpha + \beta < 1,
\]
where $K$ is physical capital, $H$ is human capital, $A$ is labour-augmenting technology,
and $L$ is labour.  Accumulation follows
\[
    \dot{K}_t = s_K Y_t - \delta K_t,\qquad
    \dot{H}_t = s_H Y_t - \delta H_t,
\]
with \emph{common} depreciation rate $\delta$, constant investment rates $s_K$ and $s_H$,
technology growth $\dot{A}/A = g > 0$, and population growth $\dot{L}/L = n > 0$.

\begin{enumerate}[label=\textbf{\arabic*.}, leftmargin=*, itemsep=6pt]

    \item \textbf{(5 points)}
    Define per-effective-worker variables $\ktilde = K/(AL)$,
    $\htilde = H/(AL)$, $\ytilde = Y/(AL)$.  Show that $\ytilde = \ktilde^{\alpha}\htilde^{\beta}$
    and derive the two laws of motion $\dot{\ktilde}$ and $\dot{\htilde}$.

    \item \textbf{(7 points)}
    Set $\dot{\ktilde} = 0$ and $\dot{\htilde} = 0$ simultaneously.
    \begin{enumerate}[label=(\alph*), itemsep=4pt]
        \item Show that the steady-state ratio satisfies $\htilde^*/\ktilde^* = s_H/s_K$.
        \item Solve for $\ktilde^*$, $\htilde^*$, and $\ytilde^*$ in closed form as
        functions of $s_K$, $s_H$, $n$, $g$, $\delta$, $\alpha$, $\beta$.
    \end{enumerate}

    \item \textbf{(7 points)}
    Take logs of the steady-state output per worker $y^* = A_t\,\ytilde^*$ to obtain the
    \textbf{MRW regression equation}:
    \[
        \ln y_t^* = \ln A_0 + g\,t
        + \frac{\alpha}{1-\alpha-\beta}\,\ln s_K
        + \frac{\beta}{1-\alpha-\beta}\,\ln s_H
        - \frac{\alpha+\beta}{1-\alpha-\beta}\,\ln(n+g+\delta).
    \]
    \begin{enumerate}[label=(\alph*), itemsep=4pt]
        \item Compare the elasticity of $y^*$ with respect to $s_K$ in this model versus
        the basic Solow model ($\beta=0$).  Evaluate both elasticities numerically for
        $\alpha = \beta = 1/3$.
        \item Explain why the augmented model better accounts for cross-country income
        differences than the basic Solow model.
    \end{enumerate}

    \item \textbf{(6 points)}
    Discuss \emph{one} important limitation of the MRW framework as a theory of
    cross-country income differences.

    \begin{hintbox}
        Think about whether $s_H$ can genuinely be treated as exogenous, and about
        what ultimately determines human capital investment.
    \end{hintbox}

\end{enumerate}

\clearpage

% ─────────────────────────────────────────────────────────────────────────────
\section{Question 2 --- Ramsey Model: Permanent Increase in TFP Growth
\hfill (25 points)}
% ─────────────────────────────────────────────────────────────────────────────

Consider a standard Ramsey/Cass/Koopmans economy.  The household has CRRA utility with
risk-aversion $\sigma > 0$ and discount rate $\rho > 0$.  Production is Cobb-Douglas:
$Y = K^{\alpha}(AL)^{1-\alpha}$.  Technology grows at rate $\dot{A}/A = x > 0$ and
population at rate $n$.  All markets are competitive.

\begin{enumerate}[label=\textbf{\arabic*.}, leftmargin=*, itemsep=6pt]

    \item \textbf{(5 points)}
    Working in per-effective-worker terms ($\ktilde = K/(AL)$,
    $\ctilde = C/(AL)$, $f(\ktilde) = \ktilde^{\alpha}$), write:
    \begin{enumerate}[label=(\alph*), itemsep=4pt]
        \item The law of motion for $\ktilde$.
        \item The Euler equation for $\ctilde$.
        \item The $\dot{\ktilde}=0$ and $\dot{\ctilde}=0$ loci, and draw the
        \textbf{phase diagram} with the saddle path for the initial economy.
    \end{enumerate}

    \item \textbf{(7 points)}
    Characterise the initial BGP:
    \begin{enumerate}[label=(\alph*), itemsep=4pt]
        \item Derive $\ktilde^*$ and $\ctilde^*$ in closed form using
        $f(\ktilde)=\ktilde^{\alpha}$.
        \item State the \textbf{modified golden rule}.
        \item State the Transversality Condition and the parameter restriction it
        imposes.
    \end{enumerate}

    \item \textbf{(7 points)}
    Suppose the economy begins at its BGP, and at $t=0$ the rate of technological
    progress \emph{permanently} increases from $x$ to $x' > x$.
    \begin{enumerate}[label=(\alph*), itemsep=4pt]
        \item Compute the \emph{new} steady state $\ktilde^{*\prime}$ and
        $\ctilde^{*\prime}$.  Show that $\ktilde^{*\prime} < \ktilde^*$ and give economic
        intuition.
        \item On the phase diagram, show how \emph{both} the $\dot{\ctilde}=0$ and
        $\dot{\ktilde}=0$ loci shift.  Clearly mark the old and new steady states.
        \item Does $\ktilde$ jump at $t=0$?  Does $\ctilde$ jump at $t=0$?  Describe the
        \textbf{transition dynamics} from the old to the new steady state: in which
        direction do $\ktilde$ and $\ctilde$ move along the new saddle path?
    \end{enumerate}

    \item \textbf{(6 points)}
    Despite the fall in per-effective-worker capital and consumption, argue that the
    household is \textbf{better off} after the shock.

    \begin{hintbox}
        Consider what happens to per-capita variables $c_t = \ctilde_t A_t$ and
        $y_t = \ytilde_t A_t$ in the long run.  Compare their growth rates before and
        after the shock.
    \end{hintbox}

\end{enumerate}

\clearpage

% ─────────────────────────────────────────────────────────────────────────────
\section{Question 3 --- Development Accounting and Misallocation
\hfill (25 points)}
% ─────────────────────────────────────────────────────────────────────────────

Consider economies with a Cobb-Douglas aggregate production function
\[
    Y = A\,K^{\alpha}\,L^{1-\alpha},\qquad \alpha \in (0,1),
\]
where $A$ is total factor productivity (TFP), $K$ is capital, and $L$ is labour.

\bigskip

\textbf{Part I: Growth and Development Accounting}

\begin{enumerate}[label=\textbf{\arabic*.}, leftmargin=*, itemsep=6pt]

    \item \textbf{(5 points)}
    \begin{enumerate}[label=(\alph*), itemsep=4pt]
        \item Using log-differentiation, derive the \textbf{growth accounting}
        decomposition:
        \[
            \gamma_Y = \gamma_A + \alpha\,\gamma_K + (1-\alpha)\,\gamma_L.
        \]
        Define the \textbf{Solow residual} and explain what it captures.
        \item List two factors besides technology that could show up in the Solow residual.
    \end{enumerate}

    \item \textbf{(7 points)}
    In per-worker terms, $y = Ak^{\alpha}$.  Consider two countries: Rich ($R$) and
    Poor ($P$) with $y_R/y_P = 27$ and $k_R/k_P = 27$.  Set $\alpha = 1/3$.
    \begin{enumerate}[label=(\alph*), itemsep=4pt]
        \item Compute the fraction of the log income gap $\ln(y_R/y_P)$ attributable to
        differences in capital per worker.
        \item Compute the fraction attributable to TFP differences.
        \item On the basis of this exercise, what is the most important proximate source
        of cross-country income differences?
    \end{enumerate}

\end{enumerate}

\bigskip

\textbf{Part II: Misallocation}

\medskip
\noindent Now consider an economy with two firms ($i=1,2$), each with one unit of labour
(fixed) and the production function $y_i = A\,k_i^{\alpha}$.  Total capital $K$ is fixed
and must be divided between the two firms: $k_1 + k_2 = K$.

\begin{enumerate}[label=\textbf{\arabic*.}, leftmargin=*, itemsep=6pt]
    \setcounter{enumi}{2}

    \item \textbf{(7 points)}
    \begin{enumerate}[label=(\alph*), itemsep=4pt]
        \item Derive the \textbf{efficient allocation} of capital.  What condition must
        hold across firms?  What is aggregate output $Y^{\mathrm{eff}}$?
        \item Now suppose firm~1 receives a government \textbf{capital subsidy}: its
        effective rental rate is $(1-\tau)r$ while firm~2 pays $r$.  Show that the
        resulting allocation satisfies
        \[
            \frac{k_1}{k_2} = (1-\tau)^{-1/(1-\alpha)}.
        \]
        \item Using the \textbf{strict concavity} of $k^{\alpha}$ (or Jensen's inequality),
        show that $Y^{\mathrm{mis}} < Y^{\mathrm{eff}}$.  Provide economic intuition.
    \end{enumerate}

    \item \textbf{(6 points)}
    \begin{enumerate}[label=(\alph*), itemsep=4pt]
        \item Define \textbf{misallocation} in terms of the dispersion of
        \emph{marginal revenue products of capital} (MRPK) across firms.
        Why does MRPK dispersion signal inefficiency?
        \item Hsieh and Klenow (2009) estimate that moving to US-level allocative
        efficiency would raise manufacturing TFP by 30--50\% in China and 40--60\% in
        India.  Give \textbf{two} concrete examples of government policies or market
        frictions that could generate the capital wedges behind these estimates.
    \end{enumerate}

\end{enumerate}

\clearpage

% ─────────────────────────────────────────────────────────────────────────────
\section{Short Questions \hfill (25 points)}
% ─────────────────────────────────────────────────────────────────────────────

\noindent\textit{For these questions, short and sharp answers will receive the highest marks.}

\vspace{8pt}

\begin{enumerate}[label=\textbf{\arabic*.}, leftmargin=*, itemsep=14pt]

    % ── SQ1 ──────────────────────────────────────────────────────────────────
    \item \textbf{(9 points) Golden Rule vs.\ Modified Golden Rule.}

    \begin{enumerate}[label=(\alph*), itemsep=6pt]
        \item In the \textbf{Solow model}, the Golden Rule level of capital maximises
        steady-state consumption.  Derive the condition $f'(k^*_{GR}) = \delta + n + g$
        and explain why it is independent of the saving rate.

        \item In the \textbf{Ramsey model}, the steady state satisfies the
        \emph{Modified Golden Rule}: $f'(\ktilde^*) = \delta + \rho + \sigma x$.
        Why does the Ramsey steady state generally \emph{not} coincide with the Golden
        Rule?  Under what condition on $\rho$ and $n$ would the two coincide?

        \item Can the Solow economy be dynamically inefficient (i.e.\ accumulate too
        much capital)?  Can the Ramsey economy?  Explain.
    \end{enumerate}

    % ── SQ2 ──────────────────────────────────────────────────────────────────
    \item \textbf{(8 points) The Role of Monopolistic Competition in Endogenous Growth.}

    \begin{enumerate}[label=(\alph*), itemsep=6pt]
        \item Explain why \emph{perfectly competitive} intermediate goods markets
        cannot sustain R\&D-based endogenous growth.  What is the fundamental problem
        with pricing new ideas at marginal cost?

        \item In the Romer model, intermediate firms charge a \textbf{markup}
        $p = 1/\alpha > 1$.  What is the source of this markup, and why is it necessary
        for the model to function?

        \item This markup creates a static distortion: intermediate goods are
        \emph{under-used} relative to the social optimum.  What policy instrument can
        correct this, and how should it be set?
    \end{enumerate}

    % ── SQ3 ──────────────────────────────────────────────────────────────────
    \item \textbf{(8 points) Expanding Varieties vs.\ Quality Ladders.}

    \begin{enumerate}[label=(\alph*), itemsep=6pt]
        \item In the \textbf{Romer (1990)} expanding varieties model, growth occurs
        through the introduction of \emph{new} goods.  In the \textbf{Aghion--Howitt
        (1992)} quality ladder model, growth occurs through \emph{improvements} to
        existing goods.  What is the key economic implication of this distinction for
        incumbent firms?

        \item Define \textbf{creative destruction} and explain the associated
        \textbf{business-stealing externality}.  Does this externality cause
        over-innovation or under-innovation relative to the social optimum?

        \item What is the \textbf{appropriability (consumer surplus) externality}?
        Does it push in the same direction as or the opposite direction from the
        business-stealing externality?  What does this imply for the overall efficiency
        of the market innovation rate?
    \end{enumerate}

\end{enumerate}

\vspace{20pt}
\noindent\rule{\textwidth}{0.4pt}

\begin{center}
    \textit{End of Practice Exam 2. Good luck!}
\end{center}

% ─────────────────────────────────────────────────────────────────────────────
\clearpage
\appendix
\section*{Appendix: Key Formulas Permitted in Exam}
% ─────────────────────────────────────────────────────────────────────────────

The following results may be used without re-derivation unless the question explicitly asks
you to derive them.

\subsection*{A.1 ~ Log-differentiation}
If $Z_t = X_t^a Y_t^b$, then $\gamma_Z = a\gamma_X + b\gamma_Y$.

\subsection*{A.2 ~ Ramsey / Euler Equation (standard, per-effective-worker)}
\[
    \frac{\dot{\ctilde}}{\ctilde}
    = \frac{1}{\sigma}\bigl(f'(\ktilde) - \delta - \rho - \sigma x\bigr),
    \qquad
    f'(\ktilde) = \alpha\ktilde^{\alpha-1}.
\]

\subsection*{A.3 ~ Solow per-effective-worker dynamics}
\[
    \dot{\ktilde} = s\,f(\ktilde) - (n+g+\delta)\,\ktilde.
\]

\subsection*{A.4 ~ Cobb-Douglas factor payments}
$rK = \alpha Y$, \quad $wL = (1-\alpha)Y$.

\subsection*{A.5 ~ Jensen's inequality}
For a strictly concave function $f$: \;
$f(x_1) + f(x_2) < 2\,f\!\bigl(\tfrac{x_1+x_2}{2}\bigr)$
whenever $x_1 \neq x_2$.

\end{document}
